% ============================================================
%  RESUMEN / ABSTRACT
% ============================================================

% ---- Resumen en español ------------------------------------
\chapter*{Resumen}
\addcontentsline{toc}{chapter}{Resumen}

El presente Trabajo Fin de Grado describe el diseño, implementación y validación de un
\textit{framework} de desarrollo rápido de aplicaciones web (\textit{Rapid Application
Development}, RAD) compuesto por dos capas complementarias: un \textit{backend} en Python
basado en \textbf{FastAPI} y \textbf{SQLAlchemy}, y un \textit{frontend} en
\textbf{Angular~19} que consume las definiciones del modelo de datos de forma declarativa.

El núcleo de la propuesta es una \textbf{tubería de transformación de esquemas}: a partir de
los modelos SQLAlchemy de la base de datos, el sistema genera automáticamente esquemas
JSON~Schema anotados con extensiones~\texttt{x-*}~(OpenAPI vendor extensions), que el
\textit{frontend} interpreta para construir formularios, tablas y operaciones CRUD sin
necesidad de escribir código repetitivo.

Entre las características implementadas destacan: un \textbf{sistema de plugins} por capas
que permite enriquecer el esquema sin modificar el modelo base, una \textbf{cuadrícula de
edición masiva} tipo Excel sobre AG~Grid~Community para inserciones y actualizaciones por
lotes, y un \textbf{editor visual de esquemas} basado en Angular~Formly.

El sistema se integra mediante contenedores \textbf{Docker} y permite desplegar nuevas
entidades con una configuración mínima, reduciendo significativamente el tiempo de
desarrollo de prototipos funcionales.

\bigskip
\noindent\textbf{Palabras clave:} desarrollo rápido de aplicaciones, FastAPI, Angular,
JSON~Schema, generación de código, bulk edit, plugins por capas, SQLAlchemy, Docker.

\vspace{2cm}

% ---- Abstract in English -----------------------------------
\chapter*{Abstract}
\addcontentsline{toc}{chapter}{Abstract}

This Bachelor's Thesis presents the design, implementation and validation of a
\textit{Rapid Application Development} (RAD) framework for web applications, consisting of
two complementary layers: a Python \textbf{FastAPI}/\textbf{SQLAlchemy} backend and an
\textbf{Angular~19} frontend that consumes data-model definitions declaratively.

The core of the proposal is a \textbf{schema transformation pipeline}: starting from
SQLAlchemy database models, the system automatically produces JSON~Schema documents
annotated with \texttt{x-*} vendor extensions (OpenAPI), which the frontend interprets to
render forms, tables and CRUD operations without boilerplate code.

Key features include: a \textbf{layered plugin system} that enriches schemas without
modifying base models, an \textbf{Excel-like bulk-edit grid} built on AG~Grid~Community
for batch inserts and updates, and a \textbf{visual schema editor} powered by
Angular~Formly.

The system is packaged as \textbf{Docker} containers, allowing new entities to be deployed
with minimal configuration and drastically reducing prototype development time.

\bigskip
\noindent\textbf{Keywords:} rapid application development, FastAPI, Angular, JSON Schema,
code generation, bulk edit, layered plugins, SQLAlchemy, Docker.
